\documentclass[10pt,a4paper]{amsart}
\usepackage[margin=19mm]{geometry}
\usepackage{amsmath,amssymb,mathtools}
\usepackage[scr=rsfs,cal=euler]{mathalfa}
\usepackage{xfrac}
\usepackage[unicode,hidelinks,bookmarks=false]{hyperref}
\newcommand{\ie}{i.e.\ }
\newcommand{\cf}{cf.\ }
\newcommand{\resp}{respectively\ }
\newcommand{\Subsection}{\subsection}
\providecommand{\nobreakdash}{\nobreak\hbox{-}}
\setlength{\emergencystretch}{2em}
\multlinegap0pt
\usepackage{mathtools}
\usepackage{amssymb}
\usepackage{tikz}
\usepackage[shortlabels]{enumitem}
\newtheorem{proposition}{Proposition}
\title{Infinite Subtraction Games with Periodic Outcomes and Aperiodic SG Values}
\author{Kai Liang}
\date{Source: arXiv:2608.17654v2 (CC BY 4.0)}
\begin{document}
\maketitle

\noindent\emph{Source and licence.} Infinite Subtraction Games with Periodic Outcomes and Aperiodic SG Values. Version arXiv:2608.17654v2 (CC BY 4.0), distributed by arXiv under the Creative Commons Attribution 4.0 International licence, http://creativecommons.org/licenses/by/4.0/, as recorded in the arXiv licence metadata for this exact version and shown on the abstract page https://arxiv.org/abs/2608.17654v2. Full licence notice: \texttt{licenses/arxiv-2608.17654-CC-BY-4.0.txt}.\par\smallskip

\noindent\textit{Editorial scope.} Counted unit: the article's proposition main\_prop (source0.tex lines 221--247). The article's complete original proof of it is delivered with it (source0.tex lines 249--318, one \texttt{proof} environment of 70 lines). No mathematics is written, repaired or re-derived: the statement and the proof are the article's own lines, in the article's own notation, and no retained line is substituted or re-flowed.  No further source-local context is needed: every cross-reference of every retained line resolves inside the two delivered spans.  Nothing was omitted from any retained span, so the package carries no omission record.  No retained line contains a citation, so the wrapper carries no bibliography and no bibliography entry.  No retained line contains a figure environment, an image-inclusion command or drawing code, so no image is delivered and the package declares no asset dependency.  Editorial formatting only, in addition: an AMS \texttt{amsart} preamble that restates the article's own declarations and macros used by the retained lines, namely the environments \texttt{proposition} on separate counters, together with the standard packages those lines need.  The retained environments print in this document as Proposition 1 in order of appearance, whereas the article prints the same items with the numbering of its own full document.  The complete native source of the article as distributed in the arXiv e-print for this exact version is bundled unchanged as \texttt{sources/207918/source0.tex}.
\begin{proposition}\label{main_prop}
    Let
    \[
        S^{>10} \subseteq (3\mathbb{N}+1)^{>10}
    \]
    be arbitrary.
    Then for the subtraction game with infinite subtraction set
    \[
        S=\{1, 4, 5, 10\} \cup S^{>10},
    \]
    the sets of positions with each SG value are
    \begin{align*}
        \mathcal{H}_0 &= \{0, 2, 8\} \cup \mathcal{H}_0{}^{>10}; \\
        \mathcal{H}_1 &= \{1, 3, 9\} \cup \mathcal{H}_1{}^{>10}; \\
        \mathcal{H}_2 &= \{4, 6\} \cup \mathcal{H}_2{}^{>10}; \\
        \mathcal{H}_3 &= \{5, 7\} \cup \mathcal{H}_3{}^{>10}; \\
        \mathcal{H}_4 &= \{10\}; \\
        \mathcal{H}_g &= \varnothing, \quad \text{for } g\geq 5,
    \end{align*}
    where
    \begin{align*}
        \mathcal{H}_0{}^{>10} &= (3\mathbb{N}+2)^{>10} = \{11, 14, 17, 20,\ldots\}; \\
        \mathcal{H}_1{}^{>10} &= (3\mathbb{N})^{>10} = \{12, 15, 18, 21,\ldots\}; \\
        \mathcal{H}_2{}^{>10} &= (3\mathbb{N}+1)^{>10} \setminus (S^{\geq 10}+6); \\
        \mathcal{H}_3{}^{>10} &= (S^{\geq 10}+6).
    \end{align*}
\end{proposition}

\begin{proof}
    By definition, to prove that a position $u\in\mathbb{N}$ has SG value $g$, it suffices to show:
    (1) for every smaller SG value $g'<g$, $u$ has at least one successor $u'$ with SG value $g'$; and
    (2) $u$ has no successor $u'$ whose SG value is still $g$.

    The cases $u\leq 10$ can be verified by direct computation, so it remains to prove the cases $g<4$ and $u>10$.

    (1) It suffices to indicate how to choose the corresponding $s$ for each $g'<g<4$.

    If $g'=0$, take
    \[
        s=
        \begin{cases}
            1, & g=1; \\
            5, & g=2,3,
        \end{cases}
    \]
    which ensures that $u-s\in \mathcal{H}_0$.
    If $g'=1$, for $g=2,3$ we simply take $s=1$.

    For the remaining case $g'=2, g=3$, we have $u\in S^{\geq 10}+6$, so taking $u'=6\in\mathcal{H}_2$ gives $s=u-u'\in S$.

    (2) That is, for every $g<4$ there are no $u,u'\in\mathcal{H}_g$ such that $u-u'\in S$.

    For $u\in\mathcal{H}_g{}^{>10}$, note that for $s\in S$,
    \[
        s\equiv
        \begin{cases}
            2, & s=5; \\
            1, & s\neq 5
        \end{cases}
        \pmod 3,
    \]
    which is never $0$; but the elements of each $\mathcal{H}_g{}^{>10}$ for $g<4$ are all congruent to a fixed residue modulo $3$, namely
    \[
        g \equiv m \coloneqq
        \begin{cases}
            2, & g=0; \\
            0, & g=1; \\
            1, & g=2,3
        \end{cases}
        \pmod 3.
    \]

    Thus $u'$ can only be one of the elements $u\in\mathcal{H}_g{}^{\leq 10}$ whose residue modulo $3$ is not $m$.
    For each $g<4$, this leaves exactly one possibility:
    \[
        u'=
        \begin{cases}
            0, & g=0; \\
            1, & g=1; \\
            6, & g=2; \\
            5, & g=3.
        \end{cases}
    \]
    Then
    \[
        s\equiv u-u' \equiv
        \begin{cases}
            2, & g=0,1,3;\\
            1, & g=2
        \end{cases}
        \pmod 3.
    \]

    Hence for $g=0,1,3$, we must have $s=5$, and the corresponding values $u=u'+s$ are $5,6,10$, none of which lie in the corresponding $\mathcal{H}_g$.

    For $g=2$, we must have $s\neq 5$; then $u=s+6$. 
    The cases $s\geq 10$ are precisely excluded by the expression for $\mathcal{H}_2$, while the remaining cases $s=1,4$ give $u=7,10$, which are not in $\mathcal{H}_2$ either.
\end{proof}

\end{document}
